This monograph presents the development of the cache memory subsystem of a RISC-V processor described in the Chisel hardware construction language, as part of a collective project whose goal is to run the Linux operating system. As the performance gap between the processor and main memory widens, the memory hierarchy ceases to be an auxiliary element and becomes a determining factor of system performance. This work specifies, implements and verifies two L1 caches, one for instructions and one for data, with parametrizable geometry, together with the miss controller, which stalls the pipeline while a miss is served, and the consistency mechanism required by the FENCE.I instruction. The objective is to quantify the effect of capacity, associativity, line size and replacement policy on the hit rate and on memory traffic, and to determine the cost in cycles of each cache event. Verification was performed with automated tests written in ScalaTest and ChiselTest, run on the Treadle simulator and compared against a reference model. The adopted configuration, with 8 KiB per cache, four ways and 16-byte lines, reached a hit rate of 69.5\% on the evaluated trace. Synthesis for the Cyclone V FPGA of the DE1-SoC board met timing at a 50 MHz clock.
